% !TEX encoding = UTF-8 Unicode
% !TEX root =  ../Main.tex

\chapter{SOLL-Analyse}\label{cha:soll}
Was ist die Idee?

Die Bearbeitung durch die Mitarbeitenden könnte auf einer gesonderten Website sstattfinden. Ein Skizze dafür ist in Abb... zu sehen. Die Idee dahinter ist, dass jede:r Mitarbeitende einen eigenen Zugang zu diesem Portal erhält. Dort werden die verschiedenen Anträge der Studierenden angezeigt. Ein Antrag ist je nach Status getagged mit "unbearbeitet, Zweitprüfung erforderlich, bewilligt, abgelehnt, archiviert". Die unbearbeiteten Anträge werden allen Mitarbeitenden angezeigt. Die Anträge, welche eine Zweitprüfung erfordern, nur den Mitarbeitenden, die den Antrag nicht erstbewilligt haben. Eine Prüfung wird gekennzeichnet, indem in der Antragsmaske ein Häkchen bei Erstpüfung und/oder Zweitprüfung gesetzt wird. Bei der Erstprüfung muss ausgewählt werden, wie die Entscheidung auf den Antrag lautet, bewilligt oder abgelehnt. Nach der Zweitprüfung erscheint dann ein Button, mit welchem der Bescheid erstellt werden kann. Dieser Bescheid wird dann automatisch an das Finanzteam sowie den/die Studierende per Mail versendet. Anträge die vollständig bearbeitet wurden oder nach bestimmter Zeit verfallen sind landen auf der Archivseite. 

    -> Vorschläge zur Verbesserung des Prozessesn und Effizienzssteigerung 
    1. Problem der Übertragung von Daten -> Lösung: Studierende tragen ihre Daten selbst in ein System ein, somit fällt dr Schritt weg, an dem wir das machen müssen -> es braucht also ein Antragsformular auf der Seite der Website
    -> Die Daten werden in einer Datenbank gespeichert inklusive notwendiger Nachweise
    -> Mitarbeitende schauen alles nochmal an, ob es irgendwelche auffälligen Fehler gibt und haken den ab 
    -> nach dem der Haken für die zweitprüfung gesetzt wurde, kann mithilfe eines Buttons eine Bewilligung erstellt werden 
    -> Bearbeitung wird viel leichter, weil der zweite Schritt die manuelle Erstellung von Bescheiden wegfällt. 
    
Ziel ist die Entwicklung einer optimierten und benutzerfreundlichen Softwarelösung zu Antragsbearbeitung. Diese Softwarelösung soll ein Formular auf der \gls{asta} Website beinhalten, mit dem Studierende einen Antrag stellen können. Dieses Formular soll ebenfalls die Option haben Dokumente hochzuladen. Da die aktuelle \gls{asta} Website mithilfe des Content-Management-Systems Wordpress erstellt wurde, soll das Antragsformular ebenfalls auf dieser Software basieren und damit funktionieren. Wenn ein:e Studierende:r einen Antrag gestellt hat, sollen die Daten sowie Dokumente in einer Datenbank gespeichert werden. Dafür sollen diese mit einer eindeutig identifizierbaren Prozessnummer gekennzeichnet werden. 



